% !TeX root = ../papers/nc-p3.tex
\section{階数15の二表示の反転と六次階数条件の縮約}
\label{sec:rank-fifteen-word-symmetries}
\begin{theorem}\label{thm:rank-fifteen-word-symmetries}
標数零の幾何点で同じ $(z,\lambda,n)$ の二表示を取る。
第二表示の道代数は第一表示の道代数の反対代数に添字 $p\mapsto-p$ で同型となる。
標準Koszul左右AS正則性の失敗候補はこの二表示で一致する。
長さ6の経路の階数低下を決めるには、第一表示の始点 $0,1$、文字 $0,2$ の
四つの行列だけを調べればよい。目標階数は文字0で36、文字2で34である。
\end{theorem}
\begin{proof}
$W$ を第一表示の四階テンソル、$W'$ を第二表示のものとする。$r^2=z,d^2=\lambda$
とし、四位置で成分順を $0,1,2,3$ とする。単元からなる次の対角係数を取る。
\[
 (d^{-1},(rd)^{-1},r^{-1},1),\quad
 (d^{-1},r/d,r^{-1},1),\quad
 (d,d/r,r^{-1},1),\quad(zd,rd,r,1).
\]
これらを $D_p$ と書くと
\[
 W'_{ijkt}=D_0(i)D_1(j)D_2(k)D_3(t)W_{t,k,j\oplus1,i\oplus1}.
\]
\path{computations/verify_rank_fifteen_word_symmetries.py} は
$(a,b,e,f)\mapsto(zb,a,zf,e)$ を元のGram表示に代入し、この恒等式を分母を払って検算する。
テンソルの各隣接切断の像が関係空間なので、一文字ごとに基底を変え語を反転すると、
第二表示の始点 $p$、長さ $m$ の商が第一表示の始点 $-p-m$ の商に移る。
文字は語の位置のうち $0,1\bmod4$ に属する個数の偶奇だけ $1$ と排他的論理和を取る。
逆写像も同じ構成で得られ、左と右の条件が交換される。

次に $\pi=(2\;3)$ を文字群の線型置換とする。同じコードは
$W_{ijkt}$ が $W_{\pi(k),\pi(t),\pi(i),\pi(j)}$ に位置ごとの対角係数を掛けたもの
であることを検算する。従って始点の二位置移動は同じパラメータでの同型で、
文字には $\pi$ が作用する。最後に正の固定窓で
$(e,f)^{\mathsf T}=u\,\mathrm{exchange}(a,b)^{\mathsf T}$ と書く
（定理\ref{thm:fixed-opposite-rank-drop}の交換行列）。
通常域ではこの行列が可逆で $u$ と $zb^2-a^2$ は単元である。
$q^2=u(zb^2-a^2)$ と $\sqrt{-1}$ を取ると、各位置で文字に順に $2,3,2,3$ を加える変換も
対角係数を掛けて $W$ を保つ。コードは二つの変換の全対角係数を与え、
全256成分について恒等式を確認する。積に現れる根の偶数乗は元の単元で評価している。
長さ6では $2,3$ が各三回現れ、全文字に $1$ が加わる。従って四始点を二つに、
四文字を二つに減らせる。経路商そのものが各局所基底で同型なので、
余核の次元を検出するFittingイデアルも対応する。単元の根を添加した後の等式は
忠実平坦な基底変換から降下する。
\end{proof}

\begin{proposition}\label{prop:rank-fifteen-oriented-cubic}
両表示の幾何学的一般窓の最小周期はちょうど2である。
同一パラメータの第一・第二表示は、一般には添字保持同型でない。
\end{proposition}
\begin{proof}
前節の $S=T^*T$ は入力側のトレース零行列空間に作用する。
$d(X,Y,Z)=\operatorname{tr}((XY+YX)Z)$ とし、トレース形式で双対な基底 $E_i,E^i$ を取る。
\[
 C=\frac{\sum_{i,j,k}d(E_i,E_j,E_k)\,d(SE^i,SE^j,SE^k)}
                {(\operatorname{tr}S)^3}.
\]
完全縮約なので基底によらず、入力の共役で $d$ は保たれる。
$T$ のスカラー倍による因子も相殺するので、これは始点を固定した有理不変量である。
前節の $x=4,c=3/4$ の101上の点で元の反対切断と一位置巡回した反対切断から
それぞれ $T$ を作り直す。上記コードによる $C$ の剰余は第一表示で $(10,44)$、
第二表示で $(44,10)$ となる。核の規格化と分母はいずれも単元で、この点は標数零に持ち上がる。
従って各表示で一位置移動との同型は一般点で不可能であり、両表示も同一パラメータでは
区別される。二位置移動との同型は前定理で構成した。
\end{proof}
